\section{Quality}

\emph{Note: This originally was part of the notes for an open class in non-duality, given in the summer of 2005.}

One of the signs of the \textbf{Kali Yuga} is the displacement of Quality by Quantity. This is so pervasive that we may find it difficult to see the distinction between the quality ``blue'' in consciousness and the quantity of a certain wavelength of light. So we have scientists since Galileo deny that the ``real'' world has any color, smell, or other sensation. Philosophers claim that ``beauty is in the eye of the beholder'' and the ``good is all relative.'' But if the scientist is not trying to manifest Truth, the artist Beauty, and the man of action Good, then what exactly are they doing?

\emph{Follow up} Since there were a few computer scientists in the group, I asked for a solution to the problem of writing a computer program to detect the color ``blue''. Every one of them suggested hooking up a spectrometer and detecting the wavelengths characterized as blue (by, of course, the programmer); they saw no paradox in this. But the true solution would be for the computer to determine the wave lengths.


\flrightit{Posted on 2010-10-24 by Cologero}

